\documentclass[10pt,a4paper]{article}
\usepackage[T1]{fontenc}
\usepackage[utf8]{inputenc}
\usepackage{lmodern}
\usepackage{amsmath,amssymb,amsthm,mathtools}
\usepackage[margin=0.88in,top=0.92in,bottom=0.88in]{geometry}
\usepackage{microtype}
\usepackage[dvipsnames]{xcolor}
\usepackage{titlesec}
\usepackage{enumitem}
\usepackage{fancyhdr}
\usepackage[hidelinks]{hyperref}

\definecolor{ink}{HTML}{17293E}
\definecolor{accent}{HTML}{275D8C}
\definecolor{muted}{HTML}{536477}
\color{ink}
\titleformat{\section}{\large\bfseries\color{accent}}{\thesection.}{0.55em}{}
\titleformat{\subsection}{\normalsize\bfseries\color{ink}}{\thesubsection}{0.55em}{}
\titlespacing*{\section}{0pt}{1.10em}{0.35em}
\titlespacing*{\subsection}{0pt}{0.75em}{0.22em}
\setlength{\parindent}{0pt}
\setlength{\parskip}{4pt}
\setlist[itemize]{leftmargin=1.35em,topsep=2pt,itemsep=1pt,parsep=0pt}
\setlist[enumerate]{leftmargin=1.5em,topsep=2pt,itemsep=1pt,parsep=0pt}
\setlength{\abovedisplayskip}{7pt}
\setlength{\belowdisplayskip}{7pt}
\setlength{\abovedisplayshortskip}{4pt}
\setlength{\belowdisplayshortskip}{4pt}
\setlength{\headheight}{20pt}
\pagestyle{fancy}
\fancyhf{}
\fancyhead[L]{\footnotesize\color{muted} STAGE 1 \quad / \quad MODULE 3}
\fancyhead[R]{\footnotesize\color{muted} MEASURE THEORY}
\fancyfoot[C]{\footnotesize\color{muted}\thepage}
\renewcommand{\headrulewidth}{0.3pt}
\renewcommand{\headrule}{\hbox to\headwidth{\color{muted!35}\leaders\hrule height \headrulewidth\hfill}}
\newcommand{\N}{\mathbb N}
\newcommand{\R}{\mathbb R}
\newcommand{\Q}{\mathbb Q}
\newcommand{\PO}{\mathcal P(\Omega)}
\newcommand{\B}{\mathcal B}
\newcommand{\key}[1]{\par\smallskip\textcolor{accent}{\textbf{Key point.}}\ #1\par}

\begin{document}
\begin{center}
  {\LARGE\bfseries Module 3: Sigma-Algebras and Measurable Spaces}\par
  \vspace{5pt}
  {\small\color{muted}Measurable structures, generators, Borel sets, and the information viewpoint}\par
  \vspace{7pt}
  {\color{accent}\rule{0.88\linewidth}{0.7pt}}
\end{center}
\vspace{-4pt}

\section{Sigma-Algebras}
\subsection{Motivation: which subsets can be measured?}
A length or probability assignment needs a domain of sets. One might try
$\mu:\mathcal P(\R)\to[0,\infty]$ with interval lengths
$\mu([a,b])=b-a$, translation invariance $\mu(A+t)=\mu(A)$,
and \emph{countable additivity} on pairwise disjoint sets. These natural requirements
cannot hold simultaneously on \emph{all} subsets of $\R$ (in the usual set-theoretic setting).

\textbf{Vitali counterexample (proof sketch).} On $[0,1]$, write
$x\sim y\iff x-y\in\Q$, and choose exactly one representative from each
equivalence class; call the resulting set $V$. This selection uses the axiom of
choice. For $q\in\Q\cap[-1,1]$, let $V_q=V+q$. By the representative property,
the sets $V_q$ are pairwise disjoint, and
\[
[0,1]\subseteq U:=\bigcup_{q\in\Q\cap[-1,1]}V_q\subseteq[-1,2].
\]
Were every subset measurable, translation invariance and countable additivity
would imply, with $a=\mu(V)$,
\[
1\le\mu(U)\le3,\qquad
\mu(U)=\sum_{q\in\Q\cap[-1,1]}\mu(V_q)=\sum_{q\in\Q\cap[-1,1]}a
=\begin{cases}0,&a=0,\\+\infty,&a>0.\end{cases}
\]
Both possibilities contradict $1\le\mu(U)\le3$. Thus a length-like measure
must restrict its domain; the Vitali set cannot be Lebesgue measurable.
The restriction should still allow complements and countable set operations.

\subsection{Definition and elementary properties}
A \emph{sigma-algebra} on $\Omega$ is a family $\mathcal F\subseteq\PO$ such that
\[
\boxed{\begin{aligned}
&\Omega\in\mathcal F,\\
&A\in\mathcal F\Longrightarrow A^c\in\mathcal F,\\
&A_n\in\mathcal F\ (n\ge1)\Longrightarrow\bigcup_{n=1}^{\infty}A_n\in\mathcal F.
\end{aligned}}
\]
Its members are called \emph{measurable sets}. By complements and De Morgan's
laws, $\varnothing\in\mathcal F$ and $\bigcap_n A_n\in\mathcal F$ whenever
$A_n\in\mathcal F$; finite unions, intersections, and differences are also measurable.
Consequently, the $\limsup$ and $\liminf$ of any sequence of measurable sets
are measurable.

Examples are the \emph{trivial} sigma-algebra $\{\varnothing,\Omega\}$ and the
\emph{discrete} sigma-algebra $\PO$. For $\Omega=\{1,2,3,4\}$,
\[
\mathcal F=\{\varnothing,\{1,2\},\{3,4\},\Omega\}
\]
is a nontrivial example: it distinguishes the two blocks but not the individual
points within either block.

\subsection{Countable, not arbitrary, additivity}
The \emph{sigma-algebra axioms} govern which unions stay in the domain;
\emph{countable additivity} is a separate axiom for a measure $\mu$ defined on
that domain:
\[
\boxed{A_i\cap A_j=\varnothing\ (i\ne j)
\quad\Longrightarrow\quad
\mu\!\left(\bigcup_{n=1}^{\infty}A_n\right)
=\sum_{n=1}^{\infty}\mu(A_n).}
\]
Arbitrary (possibly uncountable) additivity is \emph{not} required. For instance,
$[0,1]=\bigcup_{x\in[0,1]}\{x\}$, and each singleton has Lebesgue length zero,
yet $\lambda([0,1])=1$. Countable additivity permits countable limiting
operations without asserting an uncountable sum formula. Likewise, requiring
a sigma-algebra to be closed under \emph{arbitrary} unions would force it to be
$\mathcal P(\R)$ if all singletons were measurable.

\key{A sigma-algebra specifies an admissible domain for measuring sets; it does not itself assign numerical measures. Closure of the domain and additivity of a measure are different conditions.}

\section{Generated Sigma-Algebras}
\subsection{Motivation and definition}
Frequently only a collection $\mathcal C\subseteq\PO$ of basic events is given.
It need not be a sigma-algebra, so we seek the \emph{smallest} sigma-algebra
containing these events and all their logically required countable combinations:
\[
\boxed{\sigma(\mathcal C)
:=\text{the smallest sigma-algebra on $\Omega$ containing $\mathcal C$}.}
\]
This provides the minimal \emph{family closed under the sigma-algebra
operations}, not necessarily the entire power set.

\subsection{Why does the smallest one exist?}
Let $\mathfrak S=\{\mathcal F:\mathcal F\text{ is a sigma-algebra on }\Omega,
\ \mathcal C\subseteq\mathcal F\}$. It is nonempty because
$\PO\in\mathfrak S$. Define
\[
\boxed{\sigma(\mathcal C)=\bigcap_{\mathcal F\in\mathfrak S}\mathcal F.}
\]
An arbitrary intersection of sigma-algebras is again a sigma-algebra:
$\Omega$ lies in each one, and membership in \emph{all} of them is preserved
by complements and countable unions. This intersection contains $\mathcal C$
and lies inside every sigma-algebra containing $\mathcal C$, proving minimality.

For $A\subseteq\Omega$, $\sigma(\{A\})=\{\varnothing,A,A^c,\Omega\}$.
For $\Omega=\{1,2,3,4\}$, $A=\{1,2\}$ and $B=\{2,3\}$,
$\sigma(\{A,B\})=\PO$: the intersections and differences isolate all four
singletons, whose finite unions generate every subset.

\key{Generation is a minimal-closure construction: start with elementary sets and adjoin exactly what the sigma-algebra axioms force.}

\section{Borel Sigma-Algebras}
\subsection{From topology to measurability}
A topology $\mathcal T$ contains $\varnothing,X$, is closed under arbitrary
unions, and under finite intersections. It generally fails to be closed under
complements: in $\R$, $(0,1)$ is open but its complement is not.
Thus a topology need not be a sigma-algebra. The natural measurable structure
associated with a topological space $(X,\mathcal T)$ is
\[
\boxed{\B(X):=\sigma(\mathcal T),}
\]
the \emph{Borel sigma-algebra}. It is the smallest sigma-algebra containing all
open sets. In the usual topology of $\R$,
\[
\boxed{\mathcal T_{\mathrm{usual}}
\subsetneq\B(\R)\subsetneq\mathcal P(\R).}
\]
The first inclusion is strict because closed non-open sets (e.g. $[0,1]$)
are Borel; the second is strict because non-Borel subsets exist.

\subsection{Density and rational neighborhoods}
For a topological space $(X,\mathcal T)$, a subset $D\subseteq X$ is
\emph{dense} in $X$ if
\[
\boxed{\overline D=X\quad\Longleftrightarrow\quad
D\cap U\ne\varnothing\ \text{for every nonempty open }U\subseteq X.}
\]
In particular, $\Q$ is dense in $\R$. To see this, take any $a<b$.
By the \emph{Archimedean property}, choose $n\in\N$ so that
$n(b-a)>1$. Setting $k=\lfloor na\rfloor+1$ gives
\[
na<k\le na+1<nb,
\qquad\text{hence}\qquad a<\frac{k}{n}<b.
\]
Thus every nonempty open interval contains a rational number. More strongly,
for any $x\in\R$ and open neighborhood $U$ of $x$, choose
$a<x<b$ with $(a,b)\subseteq U$. Rational density provides
$p\in\Q\cap(a,x)$ and $q\in\Q\cap(x,b)$, so
\[
\boxed{x\in(p,q)\subseteq U,\qquad p,q\in\Q.}
\]
Hence intervals with rational endpoints form a \emph{countable base}
for the usual topology, motivating the countable Borel generator below.

\subsection{Examples and countable generation}
All open and closed subsets of $\R$ are Borel. Singletons $\{x\}$ are closed,
so every countable set is Borel by countable union; in particular
$\Q$ and, by complementation, $\R\setminus\Q$ are Borel. Intervals such as
$[a,b)$ are Borel as intersections of suitable closed and open sets.

A remarkably small family generates all Borel sets:
\[
\boxed{\B(\R)=\sigma\bigl(\{(p,q):p,q\in\Q,\ p<q\}\bigr).}
\]
\emph{Reason.} By the rational-neighborhood property proved above,
every open $U\subseteq\R$ is a union of rational-endpoint intervals:
\[
U=\bigcup_{\substack{p,q\in\Q,\ p<q\\(p,q)\subseteq U}}(p,q).
\]
Since $\Q\times\Q$ is countable, this is a \emph{countable} union, so
$U$ belongs to the generated sigma-algebra. Conversely, every generating
interval is open and hence Borel.

\key{Borel measurability is generated from topology. Countably many rational intervals already encode every open set in $\R$ through countable unions.}

\section{Measurable Spaces and Information}
A \emph{measurable space} is a pair
\[
\boxed{(\Omega,\mathcal F),\qquad\mathcal F\text{ a sigma-algebra on }\Omega.}
\]
No numerical size has yet been assigned. A \emph{measure space} additionally
specifies $\mu:\mathcal F\to[0,\infty]$ satisfying the measure axioms;
probability is the special case $\mathbb P(\Omega)=1$.

A sigma-algebra also describes \emph{measurable information}. For example,
let $\Omega=\{1,2,3,4\}$ and consider
\[
\mathcal F=\{\varnothing,\{1,2\},\{3,4\},\Omega\}.
\]
This distinguishes which block occurred but cannot distinguish outcomes
$1$ from $2$.
If $\mathcal F_1\subseteq\mathcal F_2$, then $\mathcal F_2$ is \emph{finer}:
it allows at least as many measurable distinctions.

For an observation $X:\Omega\to\R$, define the information it generates by
\[
\boxed{\sigma(X):=\{X^{-1}(B):B\in\B(\R)\}.}
\]
The inverse-image identities guarantee that this is a sigma-algebra. If
$X(1)=X(2)=0$ and $X(3)=X(4)=1$, then
$\sigma(X)=\{\varnothing,\{1,2\},\{3,4\},\Omega\}$.
This is the prototype for information carried by random variables and filtrations.

\section{Topology, Sigma-Algebras, and Maps}
Both topologies and sigma-algebras are families of subsets stable under
specified operations, but their closure requirements differ:
\begin{itemize}
\item \textbf{Topology:} arbitrary unions and finite intersections; complements need not be open.
\item \textbf{Sigma-algebra:} complements and countable unions (hence countable intersections); arbitrary unions need not be measurable.
\end{itemize}
They support parallel notions of structure-preserving maps.
For topological spaces $(X,\mathcal T_X)$, $(Y,\mathcal T_Y)$,
\[
\boxed{f:X\to Y\text{ continuous}
\iff f^{-1}(U)\in\mathcal T_X\quad\forall U\in\mathcal T_Y.}
\]
For measurable spaces $(X,\mathcal F)$, $(Y,\mathcal G)$,
\[
\boxed{f:X\to Y\text{ measurable}
\iff f^{-1}(B)\in\mathcal F\quad\forall B\in\mathcal G.}
\]
Both definitions use \emph{preimages} because they preserve complements and
arbitrary unions/intersections. In particular, every continuous map between
topological spaces is Borel measurable when both spaces carry their Borel
sigma-algebras: the preimages of target open sets are source open sets, and
the same conclusion extends to their generated sigma-algebras.

\key{Topology encodes local closeness; a sigma-algebra encodes measurable events. Preimages provide the common language for maps preserving either structure.}

\end{document}
